\begin{afterword}
\summaryhead

The introduction explains bias with an urn of seventy white and thirty red balls. A random sample of ten may miss the true mix by chance, but such errors are right on average and shrink as the sample grows. If the white balls are heavier and sink, the sample leans one way, and more data may not help. A cognitive bias is the same kind of systematic error in thinking, and it is harder to fix because the faulty instrument is oneself.

Examples follow. People guess that a shy stranger is a librarian, though salespeople are far more common; people stay committed to sunk costs. Studies show that knowing about a bias does not protect people from it, and that people see bias in others more easily than in themselves. Biases can still be reduced, for example by thinking in frequencies. The book's approach is to teach why good reasoning works, which the introduction compares to Serfas's finding that accounting courses reduced sunk-cost bias. A last section describes the book: essays by Yudkowsky, organized around Korzybski's ``the map is not the territory.''

\responsehead

The introduction explains statistical bias well. With one urn, and one change to it (the white balls are heavier and sink), it shows the difference between random error and systematic error. It then presents cognitive bias as an analogy to the sinking balls.

The introduction then states the strongest objection to its own book. It reports, with three studies, that knowing about a bias does not make you immune to it, and that people do not see their own biases. The book's method is to teach people about biases.

The answer has two parts, and neither meets the objection. The first is a claim that ``it is possible to recognize and overcome biases,'' supported by one sentence about frequencies with no source. The second compares the book's approach, teaching ``why good reasoning works,'' with Serfas's finding that accounting courses reduced sunk-cost bias where work experience did not. The comparison is conditional (``To the extent this volume does its job''), and nothing in the introduction shows that the condition is met. Accounting courses are formal instruction with exercises; the book is a collection of blog essays. No evidence is offered that reading the essays has the effect the courses had.

The Preface, a few pages earlier, names this very approach as the author's largest mistake: not seeing that the problem ``was figuring out how to practice it, not knowing the theory.'' The introduction then describes the book's method as teaching understanding. It even quotes Serfas saying that expertise also needs the ability to spot a bias in its setting and the tools to counteract it, and it does not say how the book supplies either.

In short: a clear account of bias, followed by the strongest objection to the book, that knowing about biases does not remove them, which the introduction answers with one unsourced sentence and a conditional comparison to accounting courses.
\end{afterword}
